\IfFileExists{stacks-project.cls}{%
\documentclass{stacks-project}
}{%
\documentclass{amsart}
}

% For dealing with references we use the comment environment
\usepackage{verbatim}
\newenvironment{reference}{\comment}{\endcomment}
%\newenvironment{reference}{}{}
\newenvironment{slogan}{\comment}{\endcomment}
\newenvironment{history}{\comment}{\endcomment}

% For commutative diagrams we use Xy-pic
\usepackage[all]{xy}

% We use 2cell for 2-commutative diagrams.
\xyoption{2cell}
\UseAllTwocells

% We use multicol for the list of chapters between chapters
\usepackage{multicol}

% This is generally recommended for better output
\usepackage{lmodern}
\usepackage[T1]{fontenc}

% For cross-file-references
\usepackage{xr-hyper}

% Package for hypertext links:
\usepackage{hyperref}

% For any local file, say "hello.tex" you want to link to please
% use \externaldocument[hello-]{hello}
\externaldocument[introduction-]{introduction}
\externaldocument[conventions-]{conventions}
\externaldocument[sets-]{sets}
\externaldocument[categories-]{categories}
\externaldocument[topology-]{topology}
\externaldocument[sheaves-]{sheaves}
\externaldocument[sites-]{sites}
\externaldocument[stacks-]{stacks}
\externaldocument[fields-]{fields}
\externaldocument[algebra-]{algebra}
\externaldocument[brauer-]{brauer}
\externaldocument[homology-]{homology}
\externaldocument[derived-]{derived}
\externaldocument[simplicial-]{simplicial}
\externaldocument[more-algebra-]{more-algebra}
\externaldocument[smoothing-]{smoothing}
\externaldocument[modules-]{modules}
\externaldocument[sites-modules-]{sites-modules}
\externaldocument[injectives-]{injectives}
\externaldocument[cohomology-]{cohomology}
\externaldocument[sites-cohomology-]{sites-cohomology}
\externaldocument[dga-]{dga}
\externaldocument[dpa-]{dpa}
\externaldocument[sdga-]{sdga}
\externaldocument[hypercovering-]{hypercovering}
\externaldocument[schemes-]{schemes}
\externaldocument[constructions-]{constructions}
\externaldocument[properties-]{properties}
\externaldocument[morphisms-]{morphisms}
\externaldocument[coherent-]{coherent}
\externaldocument[divisors-]{divisors}
\externaldocument[limits-]{limits}
\externaldocument[varieties-]{varieties}
\externaldocument[topologies-]{topologies}
\externaldocument[descent-]{descent}
\externaldocument[perfect-]{perfect}
\externaldocument[more-morphisms-]{more-morphisms}
\externaldocument[flat-]{flat}
\externaldocument[groupoids-]{groupoids}
\externaldocument[more-groupoids-]{more-groupoids}
\externaldocument[etale-]{etale}
\externaldocument[chow-]{chow}
\externaldocument[intersection-]{intersection}
\externaldocument[pic-]{pic}
\externaldocument[weil-]{weil}
\externaldocument[adequate-]{adequate}
\externaldocument[dualizing-]{dualizing}
\externaldocument[duality-]{duality}
\externaldocument[discriminant-]{discriminant}
\externaldocument[derham-]{derham}
\externaldocument[local-cohomology-]{local-cohomology}
\externaldocument[algebraization-]{algebraization}
\externaldocument[curves-]{curves}
\externaldocument[resolve-]{resolve}
\externaldocument[models-]{models}
\externaldocument[functors-]{functors}
\externaldocument[equiv-]{equiv}
\externaldocument[pione-]{pione}
\externaldocument[etale-cohomology-]{etale-cohomology}
\externaldocument[proetale-]{proetale}
\externaldocument[relative-cycles-]{relative-cycles}
\externaldocument[more-etale-]{more-etale}
\externaldocument[trace-]{trace}
\externaldocument[crystalline-]{crystalline}
\externaldocument[spaces-]{spaces}
\externaldocument[spaces-properties-]{spaces-properties}
\externaldocument[spaces-morphisms-]{spaces-morphisms}
\externaldocument[decent-spaces-]{decent-spaces}
\externaldocument[spaces-cohomology-]{spaces-cohomology}
\externaldocument[spaces-limits-]{spaces-limits}
\externaldocument[spaces-divisors-]{spaces-divisors}
\externaldocument[spaces-over-fields-]{spaces-over-fields}
\externaldocument[spaces-topologies-]{spaces-topologies}
\externaldocument[spaces-descent-]{spaces-descent}
\externaldocument[spaces-perfect-]{spaces-perfect}
\externaldocument[spaces-more-morphisms-]{spaces-more-morphisms}
\externaldocument[spaces-flat-]{spaces-flat}
\externaldocument[spaces-groupoids-]{spaces-groupoids}
\externaldocument[spaces-more-groupoids-]{spaces-more-groupoids}
\externaldocument[bootstrap-]{bootstrap}
\externaldocument[spaces-pushouts-]{spaces-pushouts}
\externaldocument[spaces-chow-]{spaces-chow}
\externaldocument[groupoids-quotients-]{groupoids-quotients}
\externaldocument[spaces-more-cohomology-]{spaces-more-cohomology}
\externaldocument[spaces-simplicial-]{spaces-simplicial}
\externaldocument[spaces-duality-]{spaces-duality}
\externaldocument[formal-spaces-]{formal-spaces}
\externaldocument[restricted-]{restricted}
\externaldocument[spaces-resolve-]{spaces-resolve}
\externaldocument[formal-defos-]{formal-defos}
\externaldocument[defos-]{defos}
\externaldocument[cotangent-]{cotangent}
\externaldocument[examples-defos-]{examples-defos}
\externaldocument[algebraic-]{algebraic}
\externaldocument[examples-stacks-]{examples-stacks}
\externaldocument[stacks-sheaves-]{stacks-sheaves}
\externaldocument[criteria-]{criteria}
\externaldocument[artin-]{artin}
\externaldocument[quot-]{quot}
\externaldocument[stacks-properties-]{stacks-properties}
\externaldocument[stacks-morphisms-]{stacks-morphisms}
\externaldocument[stacks-limits-]{stacks-limits}
\externaldocument[stacks-cohomology-]{stacks-cohomology}
\externaldocument[stacks-perfect-]{stacks-perfect}
\externaldocument[stacks-introduction-]{stacks-introduction}
\externaldocument[stacks-more-morphisms-]{stacks-more-morphisms}
\externaldocument[stacks-geometry-]{stacks-geometry}
\externaldocument[moduli-]{moduli}
\externaldocument[moduli-curves-]{moduli-curves}
\externaldocument[examples-]{examples}
\externaldocument[exercises-]{exercises}
\externaldocument[guide-]{guide}
\externaldocument[desirables-]{desirables}
\externaldocument[coding-]{coding}
\externaldocument[obsolete-]{obsolete}
\externaldocument[fdl-]{fdl}
\externaldocument[index-]{index}

% Theorem environments.
%
\theoremstyle{plain}
\newtheorem{theorem}[subsection]{Theorem}
\newtheorem{proposition}[subsection]{Proposition}
\newtheorem{lemma}[subsection]{Lemma}

\theoremstyle{definition}
\newtheorem{definition}[subsection]{Definition}
\newtheorem{example}[subsection]{Example}
\newtheorem{exercise}[subsection]{Exercise}
\newtheorem{situation}[subsection]{Situation}

\theoremstyle{remark}
\newtheorem{remark}[subsection]{Remark}
\newtheorem{remarks}[subsection]{Remarks}

\numberwithin{equation}{subsection}

% Macros
%
\def\lim{\mathop{\mathrm{lim}}\nolimits}
\def\colim{\mathop{\mathrm{colim}}\nolimits}
\def\Spec{\mathop{\mathrm{Spec}}}
\def\Hom{\mathop{\mathrm{Hom}}\nolimits}
\def\Ext{\mathop{\mathrm{Ext}}\nolimits}
\def\SheafHom{\mathop{\mathcal{H}\!\mathit{om}}\nolimits}
\def\SheafExt{\mathop{\mathcal{E}\!\mathit{xt}}\nolimits}
\def\Sch{\mathit{Sch}}
\def\Mor{\mathop{\mathrm{Mor}}\nolimits}
\def\Ob{\mathop{\mathrm{Ob}}\nolimits}
\def\Sh{\mathop{\mathit{Sh}}\nolimits}
\def\NL{\mathop{N\!L}\nolimits}
\def\CH{\mathop{\mathrm{CH}}\nolimits}
\def\proetale{{pro\text{-}\acute{e}tale}}
\def\etale{{\acute{e}tale}}
\def\QCoh{\mathit{QCoh}}
\def\Ker{\mathop{\mathrm{Ker}}}
\def\Im{\mathop{\mathrm{Im}}}
\def\Coker{\mathop{\mathrm{Coker}}}
\def\Coim{\mathop{\mathrm{Coim}}}

% Boxtimes
%
\DeclareMathSymbol{\boxtimes}{\mathbin}{AMSa}{"02}

%
% Macros for moduli stacks/spaces
%
\def\QCohstack{\mathcal{QC}\!\mathit{oh}}
\def\Cohstack{\mathcal{C}\!\mathit{oh}}
\def\Spacesstack{\mathcal{S}\!\mathit{paces}}
\def\Quotfunctor{\mathrm{Quot}}
\def\Hilbfunctor{\mathrm{Hilb}}
\def\Curvesstack{\mathcal{C}\!\mathit{urves}}
\def\Polarizedstack{\mathcal{P}\!\mathit{olarized}}
\def\Complexesstack{\mathcal{C}\!\mathit{omplexes}}
% \Pic is the operator that assigns to X its picard group, usage \Pic(X)
% \Picardstack_{X/B} denotes the Picard stack of X over B
% \Picardfunctor_{X/B} denotes the Picard functor of X over B
\def\Pic{\mathop{\mathrm{Pic}}\nolimits}
\def\Picardstack{\mathcal{P}\!\mathit{ic}}
\def\Picardfunctor{\mathrm{Pic}}
\def\Deformationcategory{\mathcal{D}\!\mathit{ef}}
\usepackage{cleveref}
\crefname{originallemma}{original lemma}{original lemmas}
\Crefname{originallemma}{Original lemma}{Original lemmas}
\makeatletter
\newlabel{lemma-2-product-fibred-categories-over-C}{{02XQ}{source}{}{sourcepointer-1}{}}
\newlabel{lemma-2-product-fibred-categories-over-C@cref}{{[originallemma][0][]02XQ}{[1][0][]source}{}{}{}}
\makeatother
\crefname{originallemma}{original lemma}{original lemmas}
\Crefname{originallemma}{Original lemma}{Original lemmas}
\makeatletter
\newlabel{lemma-fibred-equivalent}{{042G}{source}{}{sourcepointer-2}{}}
\newlabel{lemma-fibred-equivalent@cref}{{[originallemma][0][]042G}{[1][0][]source}{}{}{}}
\makeatother
\title{The inertia fibred category is equivalent to the diagonal self product}
\author{The Stacks Project Authors}
\date{stacks-project:categories.tex:local-snapshot-20261012; GNU FDL 1.2 or later}
\begin{document}
\maketitle
\noindent\textit{Editorial note.} Counted span is the author's own proof environment, source lines 7036-7093, following the statement at 6996-7034.\par\medskip
\noindent\textit{Editorial note.} Author writes 'Some details omitted' at 7092 covering well-definedness of the quasi-inverse on morphisms and naturality of the displayed isomorphism; both are routine checks of data displayed in-span.\par\medskip
\noindent\textit{Editorial note.} Steps the author leaves to the reader are flagged here for the reviewer: source line(s) 7092 carry an "omitted"/"follows easily"/"immediate from the definitions" phrase. Those sentences are retained verbatim; the displayed argument around them stays inside the counted span.\par\medskip
\noindent\textit{Allowed standard prerequisites (editorial).}\begin{itemize}
\item 2-fibre product of fibred categories over C. Lemma lemma-2-product-fibred-categories-over-C of this chapter, invoked at proof lines 7037-7042. The explicit construction of S x\_\{S'\} S and of the diagonal self product whose objects are triples (x, y, (iota, kappa)).
\item Fibred equivalence invariance. Lemma lemma-fibred-equivalent of this chapter, invoked at proof lines 7038-7039. Deduces (2) (fibredness of the inertia category) from (1) (the equivalence).
\end{itemize}\medskip
\ifdefined\hypertarget\hypertarget{sourcepointer-1}{}\fi
\noindent\textit{Source reference.} Original lemma 02XQ: sources/stacks-categories/pinned-original-source.tex, source lines 6775--6775. Navigation pointer only: the counted proof cites this statement by \textbackslash{}ref as a standing input; the deduction the counted proof has to make is carried out inside the counted span, and no part of that deduction is delegated to this pointer.\par\medskip
\ifdefined\hypertarget\hypertarget{sourcepointer-2}{}\fi
\noindent\textit{Source reference.} Original lemma 042G: sources/stacks-categories/pinned-original-source.tex, source lines 6696--6696. Navigation pointer only: the counted proof cites this statement by \textbackslash{}ref as a standing input; the deduction the counted proof has to make is carried out inside the counted span, and no part of that deduction is delegated to this pointer.\par\medskip
\begin{lemma}
\label{lemma-inertia-fibred-category}
Let $\mathcal{C}$ be a category. Let
$p : \mathcal{S} \to \mathcal{C}$ and
$p' : \mathcal{S}' \to \mathcal{C}$ be fibred categories.
Let $F : \mathcal{S} \to \mathcal{S}'$ be a $1$-morphism of
fibred categories over $\mathcal{C}$. Consider the category
$\mathcal{I}_{\mathcal{S}/\mathcal{S}'}$ over $\mathcal{C}$ whose
\begin{enumerate}
\item objects are pairs $(x, \alpha)$ where $x \in \Ob(\mathcal{S})$
and $\alpha : x \to x$ is an automorphism with $F(\alpha) = \text{id}$,
\item morphisms $(x, \alpha) \to (y, \beta)$ are given by morphisms
$\phi : x \to y$ such that
$$
\xymatrix{
x\ar[r]_\phi\ar[d]_\alpha &
y\ar[d]^{\beta} \\
x\ar[r]^\phi &
y \\
}
$$
commutes, and
\item the functor $\mathcal{I}_{\mathcal{S}/\mathcal{S}'} \to \mathcal{C}$
is given by $(x, \alpha) \mapsto p(x)$.
\end{enumerate}
Then
\begin{enumerate}
\item there is an equivalence
$$
\mathcal{I}_{\mathcal{S}/\mathcal{S}'} \longrightarrow
\mathcal{S}
\times_{\Delta, (\mathcal{S} \times_{\mathcal{S}'} \mathcal{S}), \Delta}
\mathcal{S}
$$
in the $(2, 1)$-category of categories over $\mathcal{C}$, and
\item $\mathcal{I}_{\mathcal{S}/\mathcal{S}'}$ is a fibred category over
$\mathcal{C}$.
\end{enumerate}
\end{lemma}

\begin{proof}
Note that (2) follows from (1) by
Lemmas \ref{lemma-2-product-fibred-categories-over-C} and
\ref{lemma-fibred-equivalent}. Thus it suffices to prove (1).
We will use without further mention the construction of the $2$-fibre product
from
Lemma \ref{lemma-2-product-fibred-categories-over-C}.
In particular an object of
$\mathcal{S}
\times_{\Delta, (\mathcal{S} \times_{\mathcal{S}'} \mathcal{S}), \Delta}
\mathcal{S}$
is a triple $(x, y, (\iota, \kappa))$ where $x$ and $y$ are objects of
$\mathcal{S}$, and
$(\iota, \kappa) : (x, x, \text{id}_{F(x)}) \to (y, y, \text{id}_{F(y)})$
is an isomorphism in $\mathcal{S} \times_{\mathcal{S}'} \mathcal{S}$.
This just means that $\iota, \kappa : x \to y$ are isomorphisms and that
$F(\iota) = F(\kappa)$. Consider the functor
$$
I_{\mathcal{S}/\mathcal{S}'}
\longrightarrow
\mathcal{S}
\times_{\Delta, (\mathcal{S} \times_{\mathcal{S}'} \mathcal{S}), \Delta}
\mathcal{S}
$$
which to an object $(x, \alpha)$ of the left hand side assigns the object
$(x, x, (\alpha, \text{id}_x))$ of the right hand side
and to a morphism $\phi$ of the left hand side
assigns the morphism $(\phi, \phi)$ of the right hand side.
We claim that a quasi-inverse to that morphism is given by the
functor
$$
\mathcal{S}
\times_{\Delta, (\mathcal{S} \times_{\mathcal{S}'} \mathcal{S}), \Delta}
\mathcal{S}
\longrightarrow
I_{\mathcal{S}/\mathcal{S}'}
$$
which to an object $(x, y, (\iota, \kappa))$ of the left hand side
assigns the object $(x, \kappa^{-1} \circ \iota)$ of the right hand side
and to a morphism
$(\phi, \phi') : (x, y, (\iota, \kappa)) \to (z, w, (\lambda, \mu))$
of the left hand side assigns the morphism $\phi$.
Indeed, the endo-functor of $I_{\mathcal{S}/\mathcal{S}'}$ induced
by composing the two functors above is the identity on the nose, and
the endo-functor induced on
$\mathcal{S}
\times_{\Delta, (\mathcal{S} \times_{\mathcal{S}'} \mathcal{S}), \Delta}
\mathcal{S}$
is isomorphic to
the identity via the natural isomorphism
$$
(\text{id}_x, \kappa) :
(x, x, (\kappa^{-1} \circ \iota, \text{id}_x))
\longrightarrow
(x, y, (\iota, \kappa)).
$$
Some details omitted.
\end{proof}

\end{document}
